\section{第五章 \quad 传热分析基础}

%\subsection{传热学概论}

传热学是研究在温差作用下热量传递规律的一门学科

主要研究：热量传递机理：热能如何通过导热、对流、辐射传递。热量传递规律：传递热量与温差之间的定量关系，以及温度分布的数学描述。热量传递规律的测试方法：通过实验测试获得传热过程规律

传热学研究的三类问题：强化传热、传热削弱和温度控制。

传热学的热阻思想：$\Phi =\lambda A\frac{\Delta t}{\delta}=\frac{\Delta t}{\delta / (\lambda A)}$ 我们称 $\Delta t$ 为温压，$\delta . (\lambda/A)$ 为热阻 $R_\lambda$，热阻可以串联、并联等

三种基本传热方式：热传导、热对流和热辐射

\pt{热传导}是温度不同的物体各部分之间，或温度不同且直接接触的两物体之间，依靠分子、原子、自由电子等微观粒子热运动进行的热量传递。导热是物体的固有本质，只要存在温差，固体、液体、气体中都可以发生导热。热传导满足\pt{傅里叶定律}：$q=\frac{\Phi}{A}=-\lambda \frac{\mathrm{d}t}{\mathrm{d}x}$。$\Phi$ 为热流量，单位 $\mathrm{W}$；$q$ 为热流密度，单位 $\mathrm{W}/\mathrm{m}^2$；$\lambda$ 为导热系数，单位 $\mathrm{W}/(\mathrm{m}\cdot\mathrm{K})$

\pt{热对流}：流体各部分之间发生相对位移时，冷热流体相互掺混所引起的热量传递过程。只能发生在流体中；是流体宏观运动引起的能量传递与流体分子导热共同作用的结果。对流换热是流体流过温度不同的固体壁面时，流体与固体壁面之间的热量传递过程。热对流满足\pt{牛顿冷却公式} $t_w>t_f,q=\frac{\Phi}{A}=h(t_w-t_f)$，$t_f>t_w,q=\frac{\Phi}{A}=h(t_f-t_w)$。其中 $t_w$ 是壁面温度，$t_f$ 是流体温度，$A$ 为与流体接触的壁面面积，$h$ 是对流换热系数，单位 $\mathrm{W}/(\mathrm{m}^2\cdot\mathrm{K})$。$h$ 表征对流换热强弱，是过程量，与流体种类、表面形状、流体速度等因素有关

\pt{热辐射}：物体通过电磁波传递能量；热辐射是物体由于热的原因向外发出的辐射；辐射传热是物体之间以辐射形式进行的热量交换。黑体是能吸收投入到其表面上全部热辐射能的理想物体，其辐射能力只与温度有关。黑体辐射满足\pt{斯特藩-玻尔兹曼定律}：$E=\sigma t^4$，其中 $\sigma$ 是斯特藩-玻尔兹曼常数，数值为 $5.67\times 10^{-8}\mathrm{W}/(\mathrm{m}^2\cdot\mathrm{K}^4)$；实际物体对外辐射热流量的计算式为 $\Phi =\varepsilon A \sigma T^4$。$\varepsilon$ 为发射率/黑度，$\varepsilon\le 1$，与物体种类、表面状况、温度有关。

实际热量传递常不是单一方式：多个环节可以串联，同一环节中也可能有多种方式同时作用。\pt{传热过程}：热量由壁面一侧流体通过壁面传到另一侧流体中的过程，典型如换热器中热流体向冷流体传热.\pt{传热方程式} $\Phi=kA(t_{f1}-t_{f2})$，$k$ 为传热系数，表示整个传热过程强弱，单位 $\mathrm{W}/(\mathrm{m}^2\cdot\mathrm{K})$；$k$ 与组成过程的 $h_1$、$h_2$ 等有关。

\pt{流动传热学和工程热力学的关系}： 传热学与工程热力学都属于热学范畴；热力学第一定律是研究一切热传递现象的基本依据。两种介质分界面上的能量平衡可写为 $q_0^+=q_0^-$，本质是热力学第一定律的应用。热力学研究平衡态，更关心能量总量；传热学研究非平衡态，更关心热量传递速率，例如 $q$ 的单位为 $\mathrm{J}/(\mathrm{m}^2\cdot\mathrm{s})=\mathrm{W}/\mathrm{m}^2$。

对流换热与流体流动密不可分，流体力学是传热学的先修基础；流体力学通常研究等温流动和阻力，传热学研究有温差流动中的热量交换。

传热学的研究方法：实验研究：从传热学诞生至今一直是基本研究方法，依靠实验测定传热规律。数学理论分析：对物理过程作合理简化假设，建立数学描述，再用数学方法求解析解。计算机数值求解：把偏微分方程和边界条件在空间代表点上离散为代数方程组，再近似求解。

\divider

%\subsection{传热分析基本概念}

\pt{温度场}：各个时刻物体中各点温度值所组成的集合。按照温度场随空间与时间的变化特性，可以分为

\pt{稳态温度场} $t=f(\vec{r})$ 和 \pt{非稳态温度场} $t=f(\vec{r},t)$。按照维数可以分为 \pt{一维温度场} $t=f(x)$、\pt{二维温度场} $t=f(x,y)$ 和 \pt{三维温度场} $t=f(x,y,z)$（以及是否带时间）\\

\pt{等温线（面）}同一瞬间温度相等的各点连成的线或者面被称为等温线或者等温面。等温线或者等温面不能相交，因为同一点不可能同时有两个温度；只能在物体边界中断或完全封闭。局部热流密度与等温线垂直。若相邻等温线温差相同，等温线越密，局部热流密度越大。若等温线处处垂直于某界面，则该界面为绝热面；绝热边界满足 $\partial t/\partial n=0$。

\pt{温度梯度}：某点所在等温线与相邻等温线之间的温差与其法线距离之比取极限，是指向温度升高方向的向量。$\operatorname{grad}(t)=\lim_{n\to \infty}\left(\frac{\Delta t}{\Delta n}\right)=\frac{\partial t}{\partial n}\vec{n}=\frac{\partial t}{\partial x}\vec{i}+\frac{\partial t}{\partial y}\vec{j}+\frac{\partial t}{\partial n}\vec{n}$

\divider